% --- Archon physics formalization source begin ---
% archon:physics
% archon:covers PhyXMiniProblems/problem_phyx_mini_0280.lean
% archon:source-report reports/phyx_mini/problem_phyx_mini_0280.source.json

\chapter{Physics problem phyx\_mini\_0280}
\label{ch:PhyXMiniProblems_problem_phyx_mini_0280}

\paragraph{Problem source.}
\#\# Physical scenario

A simple harmonic oscillator consists of a block attached to a spring with \textbackslash{}(k = 200\textbackslash{},\textbackslash{}mathrm\{N/m\}\textbackslash{}). The block slides on a frictionless surface, with equilibrium point \textbackslash{}(x = 0\textbackslash{}) and amplitude \textbackslash{}(0.20\textbackslash{},\textbackslash{}mathrm\{m\}\textbackslash{}). A graph of the block's velocity \textbackslash{}(v\textbackslash{}) as a function of time \textbackslash{}(t\textbackslash{}) is shown in figure. The horizontal scale is set by \textbackslash{}(t\_s = 0.20\textbackslash{},\textbackslash{}mathrm\{s\}\textbackslash{}).

\#\# Question

What is its maximum kinetic energy?

\#\# Answer choices

- A: A: 3.5 J

- B: B: 3.8 J

- C: C: 4.5 J

- D: D: 4.0 J

\#\# Figure caption (auxiliary; use the image as primary evidence)

The image is a graph depicting a sinusoidal wave. Here's the detailed description:

- The vertical axis is labeled "v (m/s)" indicating the variable represents velocity in meters per second.
- The horizontal axis is labeled "t (s)" signifying time in seconds.
- The waveform oscillates between \textbackslash{}(2\textbackslash{}pi\textbackslash{}) and \textbackslash{}(-2\textbackslash{}pi\textbackslash{}), with these values marked on the vertical axis.
- The curve starts from zero, peaks at \textbackslash{}(2\textbackslash{}pi\textbackslash{}), goes back through zero, reaches a trough at \textbackslash{}(-2\textbackslash{}pi\textbackslash{}), and then returns to zero.
- A vertical dashed line marks a time labeled \textbackslash{}(t\_s\textbackslash{}) on the horizontal axis, indicating a specific point on the curve.
- The wave represents a complete cycle of a sine function.

The graph shows the relationship between velocity and time, with the sinusoidal pattern suggesting periodic motion.

\#\# Dataset metadata

Wave Properties; Physical Model Grounding Reasoning, Multi-Formula Reasoning

\paragraph{Recorded answer/context.}
D: D: 4.0 J

\paragraph{Figure/image path.}
/root/proposal\_for\_physic/science-mango/phyx\_mini\_run/phyx\_data/test\_image/280.png

\paragraph{Formalization target.}
create a compiling Lean file with sorry bodies at `PhyXMiniProblems/problem\_phyx\_mini\_0280.lean`. The Lean declarations must preserve the physical quantities, dimensions or dimensional roles, figure labels, governing-law hypotheses, and final relation expressed by this problem.
Use Mathlib/Physlib names found through LeanExplore where available. If a physics API is missing, introduce faithful local abstractions rather than scalar placeholder aliases.

\begin{theorem}[Physics formalization target]
\label{thm:physics:phyx_mini_0280:target}
\lean{PhyXMiniProblems.ProblemPhyXMini0280.problem_phyx_mini_0280}
\uses{def:physics:phyx-mini-0280:phyxminiproblems-problemphyxmini0280-energyinjoules, def:physics:phyx-mini-0280:phyxminiproblems-problemphyxmini0280-matchessuppliedfigure, def:physics:phyx-mini-0280:phyxminiproblems-problemphyxmini0280-springblockoscillatorsetup, def:physics:phyx-mini-0280:phyxminiproblems-problemphyxmini0280-matchesproblemstatement, def:physics:phyx-mini-0280:phyxminiproblems-problemphyxmini0280-hasphysicalparameters, def:physics:phyx-mini-0280:phyxminiproblems-problemphyxmini0280-satisfiesspringblockoscillatorlaws, def:physics:phyx-mini-0280:phyxminiproblems-problemphyxmini0280-ismaximumkineticenergy, def:physics:phyx-mini-0280:phyxminiproblems-problemphyxmini0280-recordedanswerchoice, def:physics:phyx-mini-0280:phyxminiproblems-problemphyxmini0280-matchesdisplayedenergy, def:physics:phyx-mini-0280:phyxminiproblems-problemphyxmini0280-isuniqueclosestdisplayedchoice}
The assigned autoformalize agent should translate this physics problem into Lean declarations in the covered file, with theorem and lemma proof bodies written as `by sorry`.
\end{theorem}
\begin{proof}
This is an autoformalization task, not a proof task. Produce statements that can later be proved by the physics prover without weakening the physical model.
\end{proof}

% --- Archon declaration topology begin ---
\section{Declaration topology}
\begin{definition}[Mass Quantity]
  \label{def:physics:phyx-mini-0280:phyxminiproblems-problemphyxmini0280-massquantity}
  \lean{PhyXMiniProblems.ProblemPhyXMini0280.MassQuantity}
  Physical mass, with dimension M
\end{definition}
\begin{proof}
  This is the declared model definition of Mass Quantity.
\end{proof}
\begin{definition}[Length Quantity]
  \label{def:physics:phyx-mini-0280:phyxminiproblems-problemphyxmini0280-lengthquantity}
  \lean{PhyXMiniProblems.ProblemPhyXMini0280.LengthQuantity}
  Signed one-dimensional position or displacement, with dimension L
\end{definition}
\begin{proof}
  This is the declared model definition of Length Quantity.
\end{proof}
\begin{definition}[Time Quantity]
  \label{def:physics:phyx-mini-0280:phyxminiproblems-problemphyxmini0280-timequantity}
  \lean{PhyXMiniProblems.ProblemPhyXMini0280.TimeQuantity}
  Physical duration, with dimension T
\end{definition}
\begin{proof}
  This is the declared model definition of Time Quantity.
\end{proof}
\begin{definition}[Velocity Quantity]
  \label{def:physics:phyx-mini-0280:phyxminiproblems-problemphyxmini0280-velocityquantity}
  \lean{PhyXMiniProblems.ProblemPhyXMini0280.VelocityQuantity}
  Signed one-dimensional velocity, with dimension L T⁻¹
\end{definition}
\begin{proof}
  This is the declared model definition of Velocity Quantity.
\end{proof}
\begin{definition}[Spring Constant Quantity]
  \label{def:physics:phyx-mini-0280:phyxminiproblems-problemphyxmini0280-springconstantquantity}
  \lean{PhyXMiniProblems.ProblemPhyXMini0280.SpringConstantQuantity}
  Spring stiffness, with dimension force per length, equivalently M T⁻²
\end{definition}
\begin{proof}
  This is the declared model definition of Spring Constant Quantity.
\end{proof}
\begin{definition}[Energy Quantity]
  \label{def:physics:phyx-mini-0280:phyxminiproblems-problemphyxmini0280-energyquantity}
  \lean{PhyXMiniProblems.ProblemPhyXMini0280.EnergyQuantity}
  Mechanical energy, with dimension M L² T⁻²
\end{definition}
\begin{proof}
  This is the declared model definition of Energy Quantity.
\end{proof}
\begin{definition}[mass In Kilograms]
  \label{def:physics:phyx-mini-0280:phyxminiproblems-problemphyxmini0280-massinkilograms}
  \lean{PhyXMiniProblems.ProblemPhyXMini0280.massInKilograms}
  \uses{def:physics:phyx-mini-0280:phyxminiproblems-problemphyxmini0280-massquantity}
  SI kilogram readout of a physical mass
\end{definition}
\begin{proof}
  This is the declared model definition of mass In Kilograms.
\end{proof}
\begin{definition}[length In Meters]
  \label{def:physics:phyx-mini-0280:phyxminiproblems-problemphyxmini0280-lengthinmeters}
  \lean{PhyXMiniProblems.ProblemPhyXMini0280.lengthInMeters}
  \uses{def:physics:phyx-mini-0280:phyxminiproblems-problemphyxmini0280-lengthquantity}
  SI metre readout of a signed position or displacement
\end{definition}
\begin{proof}
  This is the declared model definition of length In Meters.
\end{proof}
\begin{definition}[time In Seconds]
  \label{def:physics:phyx-mini-0280:phyxminiproblems-problemphyxmini0280-timeinseconds}
  \lean{PhyXMiniProblems.ProblemPhyXMini0280.timeInSeconds}
  \uses{def:physics:phyx-mini-0280:phyxminiproblems-problemphyxmini0280-timequantity}
  SI second readout of a physical duration
\end{definition}
\begin{proof}
  This is the declared model definition of time In Seconds.
\end{proof}
\begin{definition}[velocity In Meters Per Second]
  \label{def:physics:phyx-mini-0280:phyxminiproblems-problemphyxmini0280-velocityinmeterspersecond}
  \lean{PhyXMiniProblems.ProblemPhyXMini0280.velocityInMetersPerSecond}
  \uses{def:physics:phyx-mini-0280:phyxminiproblems-problemphyxmini0280-velocityquantity}
  SI metre-per-second readout of a signed velocity component
\end{definition}
\begin{proof}
  This is the declared model definition of velocity In Meters Per Second.
\end{proof}
\begin{definition}[spring Constant In Newtons Per Meter]
  \label{def:physics:phyx-mini-0280:phyxminiproblems-problemphyxmini0280-springconstantinnewtonspermeter}
  \lean{PhyXMiniProblems.ProblemPhyXMini0280.springConstantInNewtonsPerMeter}
  \uses{def:physics:phyx-mini-0280:phyxminiproblems-problemphyxmini0280-springconstantquantity}
  SI newton-per-metre readout of a spring stiffness
\end{definition}
\begin{proof}
  This is the declared model definition of spring Constant In Newtons Per Meter.
\end{proof}
\begin{definition}[energy In Joules]
  \label{def:physics:phyx-mini-0280:phyxminiproblems-problemphyxmini0280-energyinjoules}
  \lean{PhyXMiniProblems.ProblemPhyXMini0280.energyInJoules}
  \uses{def:physics:phyx-mini-0280:phyxminiproblems-problemphyxmini0280-energyquantity}
  SI joule readout of a physical energy
\end{definition}
\begin{proof}
  This is the declared model definition of energy In Joules.
\end{proof}
\begin{definition}[Axis Quantity]
  \label{def:physics:phyx-mini-0280:phyxminiproblems-problemphyxmini0280-axisquantity}
  \lean{PhyXMiniProblems.ProblemPhyXMini0280.AxisQuantity}
  Physical quantity assigned to a graph axis
\end{definition}
\begin{proof}
  This is the declared model definition of Axis Quantity.
\end{proof}
\begin{definition}[Axis Unit]
  \label{def:physics:phyx-mini-0280:phyxminiproblems-problemphyxmini0280-axisunit}
  \lean{PhyXMiniProblems.ProblemPhyXMini0280.AxisUnit}
  Unit printed beside a graph axis
\end{definition}
\begin{proof}
  This is the declared model definition of Axis Unit.
\end{proof}
\begin{definition}[Waveform Shape]
  \label{def:physics:phyx-mini-0280:phyxminiproblems-problemphyxmini0280-waveformshape}
  \lean{PhyXMiniProblems.ProblemPhyXMini0280.WaveformShape}
  Qualitative shape of the plotted velocity trace
\end{definition}
\begin{proof}
  This is the declared model definition of Waveform Shape.
\end{proof}
\begin{definition}[Displayed Cycle]
  \label{def:physics:phyx-mini-0280:phyxminiproblems-problemphyxmini0280-displayedcycle}
  \lean{PhyXMiniProblems.ProblemPhyXMini0280.DisplayedCycle}
  Portion of an oscillation displayed between 0 and t\_s
\end{definition}
\begin{proof}
  This is the declared model definition of Displayed Cycle.
\end{proof}
\begin{definition}[Time Marker Meaning]
  \label{def:physics:phyx-mini-0280:phyxminiproblems-problemphyxmini0280-timemarkermeaning}
  \lean{PhyXMiniProblems.ProblemPhyXMini0280.TimeMarkerMeaning}
  Meaning of the dashed vertical marker carrying the label t\_s
\end{definition}
\begin{proof}
  This is the declared model definition of Time Marker Meaning.
\end{proof}
\begin{definition}[Velocity Time Figure]
  \label{def:physics:phyx-mini-0280:phyxminiproblems-problemphyxmini0280-velocitytimefigure}
  \lean{PhyXMiniProblems.ProblemPhyXMini0280.VelocityTimeFigure}
  \uses{def:physics:phyx-mini-0280:phyxminiproblems-problemphyxmini0280-timequantity, def:physics:phyx-mini-0280:phyxminiproblems-problemphyxmini0280-velocityquantity, def:physics:phyx-mini-0280:phyxminiproblems-problemphyxmini0280-axisquantity, def:physics:phyx-mini-0280:phyxminiproblems-problemphyxmini0280-axisunit, def:physics:phyx-mini-0280:phyxminiproblems-problemphyxmini0280-waveformshape, def:physics:phyx-mini-0280:phyxminiproblems-problemphyxmini0280-displayedcycle, def:physics:phyx-mini-0280:phyxminiproblems-problemphyxmini0280-timemarkermeaning}
  The raw graph data. The argument of velocityTraceAtSeconds is the scalar time coordinate printed in seconds, while its value is a dimensionful signed velocity. The numerical extrema are stated separately in MatchesSuppliedFigure, not built into this structure
\end{definition}
\begin{proof}
  This is the declared model definition of Velocity Time Figure.
\end{proof}
\begin{definition}[Matches Supplied Figure]
  \label{def:physics:phyx-mini-0280:phyxminiproblems-problemphyxmini0280-matchessuppliedfigure}
  \lean{PhyXMiniProblems.ProblemPhyXMini0280.MatchesSuppliedFigure}
  \uses{def:physics:phyx-mini-0280:phyxminiproblems-problemphyxmini0280-timeinseconds, def:physics:phyx-mini-0280:phyxminiproblems-problemphyxmini0280-velocityinmeterspersecond, def:physics:phyx-mini-0280:phyxminiproblems-problemphyxmini0280-velocitytimefigure}
  Primary-image evidence. The velocity trace begins and ends at zero during the displayed interval, reaches both labeled extrema ±2π m/s, remains between them, and the t\_s marker denotes the end of one complete cycle
\end{definition}
\begin{proof}
  This is the declared model definition of Matches Supplied Figure.
\end{proof}
\begin{definition}[Surface Friction]
  \label{def:physics:phyx-mini-0280:phyxminiproblems-problemphyxmini0280-surfacefriction}
  \lean{PhyXMiniProblems.ProblemPhyXMini0280.SurfaceFriction}
  Contact model for the block and its support surface
\end{definition}
\begin{proof}
  This is the declared model definition of Surface Friction.
\end{proof}
\begin{definition}[Motion Regime]
  \label{def:physics:phyx-mini-0280:phyxminiproblems-problemphyxmini0280-motionregime}
  \lean{PhyXMiniProblems.ProblemPhyXMini0280.MotionRegime}
  Dynamical regime asserted in the problem statement
\end{definition}
\begin{proof}
  This is the declared model definition of Motion Regime.
\end{proof}
\begin{definition}[Spring Block Oscillator Setup]
  \label{def:physics:phyx-mini-0280:phyxminiproblems-problemphyxmini0280-springblockoscillatorsetup}
  \lean{PhyXMiniProblems.ProblemPhyXMini0280.SpringBlockOscillatorSetup}
  \uses{def:physics:phyx-mini-0280:phyxminiproblems-problemphyxmini0280-massquantity, def:physics:phyx-mini-0280:phyxminiproblems-problemphyxmini0280-lengthquantity, def:physics:phyx-mini-0280:phyxminiproblems-problemphyxmini0280-springconstantquantity, def:physics:phyx-mini-0280:phyxminiproblems-problemphyxmini0280-energyquantity, def:physics:phyx-mini-0280:phyxminiproblems-problemphyxmini0280-velocitytimefigure, def:physics:phyx-mini-0280:phyxminiproblems-problemphyxmini0280-surfacefriction, def:physics:phyx-mini-0280:phyxminiproblems-problemphyxmini0280-motionregime}
  The physical spring--block system. oscillatorSIReadout, initialConditions, and amplitudeVectorInMeters are the one-dimensional SI-coordinate data used by PhysLean. The requested maximum kinetic energy remains an unconstrained dimensionful field here; its meaning and numerical value are supplied separately
\end{definition}
\begin{proof}
  This is the declared model definition of Spring Block Oscillator Setup.
\end{proof}
\begin{definition}[Matches Problem Statement]
  \label{def:physics:phyx-mini-0280:phyxminiproblems-problemphyxmini0280-matchesproblemstatement}
  \lean{PhyXMiniProblems.ProblemPhyXMini0280.MatchesProblemStatement}
  \uses{def:physics:phyx-mini-0280:phyxminiproblems-problemphyxmini0280-lengthinmeters, def:physics:phyx-mini-0280:phyxminiproblems-problemphyxmini0280-timeinseconds, def:physics:phyx-mini-0280:phyxminiproblems-problemphyxmini0280-springconstantinnewtonspermeter, def:physics:phyx-mini-0280:phyxminiproblems-problemphyxmini0280-springblockoscillatorsetup}
  Measurements and qualitative conditions stated in the prose. Both 0.20 m and 0.20 s are represented exactly as 1/5 in SI readouts. No mass or kinetic-energy answer is assumed
\end{definition}
\begin{proof}
  This is the declared model definition of Matches Problem Statement.
\end{proof}
\begin{definition}[Has Physical Parameters]
  \label{def:physics:phyx-mini-0280:phyxminiproblems-problemphyxmini0280-hasphysicalparameters}
  \lean{PhyXMiniProblems.ProblemPhyXMini0280.HasPhysicalParameters}
  \uses{def:physics:phyx-mini-0280:phyxminiproblems-problemphyxmini0280-massinkilograms, def:physics:phyx-mini-0280:phyxminiproblems-problemphyxmini0280-lengthinmeters, def:physics:phyx-mini-0280:phyxminiproblems-problemphyxmini0280-timeinseconds, def:physics:phyx-mini-0280:phyxminiproblems-problemphyxmini0280-springconstantinnewtonspermeter, def:physics:phyx-mini-0280:phyxminiproblems-problemphyxmini0280-springblockoscillatorsetup}
  Positivity and nondegeneracy of the physical parameters
\end{definition}
\begin{proof}
  This is the declared model definition of Has Physical Parameters.
\end{proof}
\begin{definition}[Satisfies Spring Block Oscillator Laws]
  \label{def:physics:phyx-mini-0280:phyxminiproblems-problemphyxmini0280-satisfiesspringblockoscillatorlaws}
  \lean{PhyXMiniProblems.ProblemPhyXMini0280.SatisfiesSpringBlockOscillatorLaws}
  \uses{def:physics:phyx-mini-0280:phyxminiproblems-problemphyxmini0280-massinkilograms, def:physics:phyx-mini-0280:phyxminiproblems-problemphyxmini0280-lengthinmeters, def:physics:phyx-mini-0280:phyxminiproblems-problemphyxmini0280-timeinseconds, def:physics:phyx-mini-0280:phyxminiproblems-problemphyxmini0280-velocityinmeterspersecond, def:physics:phyx-mini-0280:phyxminiproblems-problemphyxmini0280-springconstantinnewtonspermeter, def:physics:phyx-mini-0280:phyxminiproblems-problemphyxmini0280-springblockoscillatorsetup}
  The governing simple-harmonic-oscillator model and the bridge from the graph to the physical trajectory. The PhysLean oscillator uses the coherent SI mass and stiffness readouts. At the graph's time origin the block is at the negative turning point and is at rest, which gives the displayed initially positive sinusoidal velocity. The graph scale is the oscillator period. The final field states the generic kinetic-energy law for the graph velocity and PhysLean trajectory;
\end{definition}
\begin{proof}
  This is the declared model definition of Satisfies Spring Block Oscillator Laws.
\end{proof}
\begin{definition}[Is Maximum Kinetic Energy]
  \label{def:physics:phyx-mini-0280:phyxminiproblems-problemphyxmini0280-ismaximumkineticenergy}
  \lean{PhyXMiniProblems.ProblemPhyXMini0280.IsMaximumKineticEnergy}
  \uses{def:physics:phyx-mini-0280:phyxminiproblems-problemphyxmini0280-energyinjoules, def:physics:phyx-mini-0280:phyxminiproblems-problemphyxmini0280-springblockoscillatorsetup}
  The named dimensionful energy is the attained maximum of the standard PhysLean kinetic-energy trace. This predicate gives the word "maximum" its mathematical meaning but does not state its value or identify an answer
\end{definition}
\begin{proof}
  This is the declared model definition of Is Maximum Kinetic Energy.
\end{proof}
\begin{lemma}[maximum Kinetic Energy eq amplitude Potential Energy]
  \label{lem:physics:phyx-mini-0280:phyxminiproblems-problemphyxmini0280-maximumkineticenergy-eq-amplitudepotentialenergy}
  \lean{PhyXMiniProblems.ProblemPhyXMini0280.maximumKineticEnergy_eq_amplitudePotentialEnergy}
  \uses{def:physics:phyx-mini-0280:phyxminiproblems-problemphyxmini0280-energyinjoules, def:physics:phyx-mini-0280:phyxminiproblems-problemphyxmini0280-springblockoscillatorsetup, def:physics:phyx-mini-0280:phyxminiproblems-problemphyxmini0280-hasphysicalparameters, def:physics:phyx-mini-0280:phyxminiproblems-problemphyxmini0280-satisfiesspringblockoscillatorlaws, def:physics:phyx-mini-0280:phyxminiproblems-problemphyxmini0280-ismaximumkineticenergy}
  For a frictionless oscillator released from a turning point, conservation of mechanical energy makes the maximum kinetic energy equal to the spring potential energy at the amplitude. This is a derived physical relation, not a field of the governing-law interface
\end{lemma}
\begin{proof}
  The conclusion follows from the governing relations and figure readouts named in the statement of maximum Kinetic Energy eq amplitude Potential Energy.
\end{proof}
\begin{definition}[Answer Choice]
  \label{def:physics:phyx-mini-0280:phyxminiproblems-problemphyxmini0280-answerchoice}
  \lean{PhyXMiniProblems.ProblemPhyXMini0280.AnswerChoice}
  Labels printed beside the four candidate energies
\end{definition}
\begin{proof}
  This is the declared model definition of Answer Choice.
\end{proof}
\begin{definition}[energy In Joules]
  \label{def:physics:phyx-mini-0280:phyxminiproblems-problemphyxmini0280-answerchoice-energyinjoules}
  \lean{PhyXMiniProblems.ProblemPhyXMini0280.AnswerChoice.energyInJoules}
  \uses{def:physics:phyx-mini-0280:phyxminiproblems-problemphyxmini0280-energyinjoules, def:physics:phyx-mini-0280:phyxminiproblems-problemphyxmini0280-answerchoice}
  Energy in joules printed beside each answer label
\end{definition}
\begin{proof}
  This is the declared model definition of energy In Joules.
\end{proof}
\begin{definition}[recorded Answer Choice]
  \label{def:physics:phyx-mini-0280:phyxminiproblems-problemphyxmini0280-recordedanswerchoice}
  \lean{PhyXMiniProblems.ProblemPhyXMini0280.recordedAnswerChoice}
  \uses{def:physics:phyx-mini-0280:phyxminiproblems-problemphyxmini0280-answerchoice}
  The answer label recorded in the supplied dataset metadata
\end{definition}
\begin{proof}
  This is the declared model definition of recorded Answer Choice.
\end{proof}
\begin{definition}[Matches Displayed Energy]
  \label{def:physics:phyx-mini-0280:phyxminiproblems-problemphyxmini0280-matchesdisplayedenergy}
  \lean{PhyXMiniProblems.ProblemPhyXMini0280.MatchesDisplayedEnergy}
  \uses{def:physics:phyx-mini-0280:phyxminiproblems-problemphyxmini0280-energyinjoules, def:physics:phyx-mini-0280:phyxminiproblems-problemphyxmini0280-springblockoscillatorsetup, def:physics:phyx-mini-0280:phyxminiproblems-problemphyxmini0280-answerchoice}
  Exact agreement of the physical energy with a displayed choice
\end{definition}
\begin{proof}
  This is the declared model definition of Matches Displayed Energy.
\end{proof}
\begin{definition}[Is Unique Closest Displayed Choice]
  \label{def:physics:phyx-mini-0280:phyxminiproblems-problemphyxmini0280-isuniqueclosestdisplayedchoice}
  \lean{PhyXMiniProblems.ProblemPhyXMini0280.IsUniqueClosestDisplayedChoice}
  \uses{def:physics:phyx-mini-0280:phyxminiproblems-problemphyxmini0280-energyinjoules, def:physics:phyx-mini-0280:phyxminiproblems-problemphyxmini0280-springblockoscillatorsetup, def:physics:phyx-mini-0280:phyxminiproblems-problemphyxmini0280-answerchoice}
  The selected displayed value is strictly closest to the physical energy
\end{definition}
\begin{proof}
  This is the declared model definition of Is Unique Closest Displayed Choice.
\end{proof}
% --- Archon declaration topology end ---
% --- Archon physics formalization source end ---
